\subsubsection{BlackBerry OS}

\paragraph{Address Translation and Memory Model}

BlackBerry~10 uses paging through the CPU's MMU. Physical memory is
divided into pages (typically 4\,KB), and hardware page tables map each
process's virtual addresses to physical addresses; the process manager
owns and updates these tables as it starts and switches processes
\citep{qnx_memman}. The translation look-aside buffer (TLB) caches
recent translations. QNX also supports variable page sizes where the
processor allows: a larger page reduces the number of TLB entries
needed, cutting TLB misses and improving memory-access speed---at the
cost of coarser allocation granularity \citep{qnx_memman}. On older ARM cores, QNX
used the FCSE slot scheme to avoid cache flushes across context
switches \citep{qnx_arm_memory}; on the ARMv7 cores in BlackBerry~10,
this role is filled by ASID-tagged TLB entries instead. The structures involved are
therefore page tables, the TLB, and per-process mappings.

\paragraph{Virtual Memory Management and Implementation}

BlackBerry~10 implements on-demand paging: when memory is mapped (for
example via \texttt{mmap()}), physical pages are generally \emph{not}
allocated immediately. Instead the mapping records what should back the range;
the first read or write to a page triggers a page fault, on which the
memory manager allocates a physical page, initialises its contents,
and installs the page-table entry \citep{membarrier_qnx8_mem}. The
performance implication is lower up-front memory use and faster
mapping, paid for by a one-time fault cost on first touch.

\paragraph{Memory Allocation Strategies}

Each process receives its own large private virtual address space (up
to 2--3.5\,GB depending on CPU) \citep{qnx_memman}. Within it, the
loader lays out the text, data, BSS, stack, and shared libraries;
\texttt{malloc()} then finds free virtual ranges for the heap
\citep{qnx_arm_memory}. The OS balances efficient allocation against
speed by allocating lazily (on demand) and by supporting POSIX memory
locking (\texttt{mlock}) so a latency-sensitive process can pin pages
resident and avoid fault latency entirely \citep{qnx_memman}---a direct
speed-versus-memory trade the designer controls.

\paragraph{Page Replacement Policies}

Page replacement matters when the system pages to backing store under
memory pressure. QNX's modern memory manager is built around on-demand
paging with a working-set discipline, choosing which pages to retain
based on usage \citep{membarrier_qnx8_mem}. As an embedded/real-time
OS, QNX is typically configured to \emph{minimise} dependence on
swapping, because the unbounded latency of fetching a page from slow
storage is hostile to real-time guarantees---hence the prominence of
memory locking to keep critical pages resident \citep{qnx_memman}. The
trade-off is accuracy (keeping the truly hot pages) versus the overhead
of tracking page usage.

\paragraph{Memory Protection}

Protection is enforced by the MMU plus user/kernel-mode separation.
Each process is isolated; page permissions govern read/write/execute
access \citep{qnx_memman}. When a process attempts an illegal access---
touching unmapped memory or violating page permissions---the MMU
notifies the OS, which aborts the offending thread at the faulting
instruction, containing the fault rather than letting it corrupt other
processes or the kernel \citep{qnx_memman}. This is the foundation of
the OS's robustness: faults are local and recoverable, which is what
allows the HAM to restart a crashed component \citep{qnx_ham}.

\paragraph{Fragmentation}

Paging itself largely eliminates \emph{external} fragmentation, since
any free physical page can back any virtual page---memory contiguous in
the program may be split across physical memory transparently
\citep{qnx_memman}. \emph{Internal} fragmentation (partly used pages)
is the residual cost, and variable page sizes are a tuning knob here
\citep{qnx_memman}. QNX also performs physical-memory defragmentation in
its memory manager rather than relying on stop-the-world compaction,
consistent with a real-time design that avoids long uninterruptible
operations \citep{qnx_memman}.

\paragraph{Behaviour Under Pressure}

Under low memory, an on-demand-paging system either reclaims/pages out
cold pages or, in a configuration without swap, fails allocations and
may terminate processes. Because real-time correctness forbids
unbounded paging latency, QNX leans toward keeping the working set
resident (memory locking) and, where the optional adaptive partitioning
scheduler is enabled, bounding usage with partitioning rather than
thrashing on a swap device \citep{qnx_memman, qnx_aps}. Thrashing
is best avoided in QNX by design (locking critical memory, budgeting
resources) rather than detected after the fact. This makes its
behaviour under pressure predictable, if less elastic than a desktop OS
that swaps aggressively.

\paragraph{Performance Considerations}

Page size, TLB efficiency, and page-fault rate are the levers. Larger
variable pages reduce TLB misses and speed access but coarsen
allocation \citep{qnx_memman}; on-demand paging trims memory at the
cost of first-touch faults \citep{membarrier_qnx8_mem}; memory locking
removes fault latency at the cost of holding physical RAM
\citep{qnx_memman}. The design choices that improve performance
(locking, larger pages) tend to increase memory usage, and vice
versa---a trade BlackBerry/QNX exposes to the system designer rather
than hiding.

\paragraph{Overview --- Memory Management}

Strengths include strong per-process isolation, predictable (lockable)
memory behaviour suited to real-time use, low external fragmentation,
and tunable page sizes \citep{qnx_memman, membarrier_qnx8_mem}.
Weaknesses are that private address spaces add per-process page-table
cost ($\approx$\,4--8\,KB each) and make cross-space operations
costlier than in a shared-space design \citep{qnx_memman}; the
real-time bias toward locking trades elasticity for predictability. On
a system with several apps and limited RAM, BlackBerry~10 keeps each
app protected and behaves predictably, but a desktop-class OS with
aggressive swapping might keep \emph{more} apps nominally alive under
heavy overcommit. For memory-intensive apps (browsers, games), its MMU
isolation and demand paging serve them well, while a designer can lock
the latency-critical portions resident \citep{qnx_memman}.
